% ==============================================================================
% PREÁMBULO UNIFICADO
% Base: proyecto original (tcolorbox, enlaces en color, márgenes de 2.5 cm).
% Incorporado del proyecto nuevo: \imagen, \huecodiagrama y los entornos
% ejercicio / resolucion / mecanizado para los exámenes resueltos.
%
% Uso: \input{preambulo} en el .tex principal (justo tras \documentclass),
% o bien pegar el contenido directamente.
% ==============================================================================

% --- Fuentes, codificación e idioma -------------------------------------------
\usepackage[T1]{fontenc}
\usepackage[utf8]{inputenc}
\usepackage{lmodern}
\usepackage[spanish, es-tabla]{babel}

% --- Matemáticas y utilidades -------------------------------------------------
\usepackage{amsmath, amssymb}   % amssymb ya carga amsfonts
\usepackage{enumitem}
\usepackage{etoolbox}           % \ifstrempty

% --- Color y cajas ------------------------------------------------------------
\usepackage{xcolor}
\usepackage[most]{tcolorbox}

% --- Figuras (del proyecto nuevo) ---------------------------------------------
\usepackage{graphicx}
\usepackage{caption}            % antes de hyperref
\graphicspath{{imagenes/}}      % si el original ya tenía otra ruta: {{imagenes/}{figuras/}}
\captionsetup{font=small,labelfont=bf}
\newcommand{\imagen}[4][0.8\linewidth]{%
  \par\addvspace{0.6em}%
  \begin{center}
    \includegraphics[width=#1]{#2}\captionof{figure}{#3}\label{#4}
  \end{center}\par\addvspace{0.6em}}

% --- Página -------------------------------------------------------------------
\usepackage[margin=2.5cm]{geometry}

% --- Hipervínculos (siempre al final de la lista de paquetes) -----------------
\usepackage{hyperref}
\hypersetup{
    colorlinks=true,
    linkcolor=blue!80!black,
    citecolor=blue!80!black,
    urlcolor=blue!80!black,
    pageanchor=false
}

% ==============================================================================
% COLORES Y ESTILOS DE CAJA (originales)
% ==============================================================================

% Paleta de colores
\definecolor{usalBlue}{RGB}{0, 51, 102}
\definecolor{usalRed}{RGB}{180, 40, 40}
\definecolor{bgBlue}{RGB}{240, 245, 255}
\definecolor{bgGreen}{RGB}{240, 250, 242}
\definecolor{bgYellow}{RGB}{255, 253, 235}

\definecolor{maincolor}{RGB}{0, 51, 102}
\definecolor{accentcolor}{RGB}{0, 128, 128}
\definecolor{boxbg}{RGB}{245, 247, 250}

% Cajas personalizadas
\tcbset{
    cajaResumen/.style={
        enhanced jigsaw, breakable,
        colback=bgBlue,
        colframe=usalBlue,
        fonttitle=\bfseries,
        coltitle=white,
        arc=2mm,
        boxrule=1pt,
        left=10pt, right=10pt, top=8pt, bottom=8pt
    },
    cajaFormulario/.style={
        enhanced jigsaw, breakable,
        colback=bgYellow,
        colframe=usalRed,
        fonttitle=\bfseries,
        coltitle=white,
        arc=2mm,
        boxrule=1pt,
        left=10pt, right=10pt, top=8pt, bottom=8pt
    }
}

% ==============================================================================
% PROBLEMAS: DOS ESTILOS CON UN MISMO NOMBRE PÚBLICO (`problema`)
% ------------------------------------------------------------------------------
% Los dos proyectos definían `problema` y `problemacontador`, así que:
%   - problemacaja : el `problema` ORIGINAL (tcolorbox, "Problema 3.1: ...")
%   - ejercicio    : el `problema` NUEVO (encabezado simple, "Ejercicio 1: ...")
% `problema` es un alias de uno u otro según el estilo activo:
%   \estiloApuntes -> `problema` = caja original   (activo por defecto)
%   \estiloExamen  -> `problema` = estilo ejercicio (para el capítulo de exámenes)
% Ambos comparten el contador `problemacontador` (se reinicia en cada capítulo).
% ==============================================================================

% Contador (reinicia en cada capítulo; con clase article usar [section])
\newcounter{problemacontador}[chapter]

% --- Estilo original: caja ----------------------------------------------------
\newtcolorbox[use counter=problemacontador]{problemacaja}[1][]{
  enhanced jigsaw,       % como enhanced, pero con esquinas rectas en los cortes
  breakable,
  colback=boxbg,
  colframe=maincolor,
  fonttitle=\bfseries,
  coltitle=white,
  arc=2mm,
  boxrule=1pt,
  before skip=10pt,
  after skip=10pt,
  title={Problema~\thetcbcounter\ifstrempty{#1}{}{: #1}},
  lower separated=true
}

% --- Estilo exámenes: encabezado simple ---------------------------------------
\newenvironment{ejercicio}[1][]{%
  \par\addvspace{1.2em}%
  % Numeración sin prefijo de capítulo ("Ejercicio 1", no "Ejercicio 5.1");
  % el cambio es local al entorno y afecta también a \label/\ref.
  \renewcommand{\theproblemacontador}{\arabic{problemacontador}}%
  \noindent\refstepcounter{problemacontador}%
  {\Large\bfseries Ejercicio \theproblemacontador
   \ifstrempty{#1}{}{: #1}\par}\nobreak
  \vspace{0.4em}\noindent\ignorespaces
}{%
  \par\addvspace{0.8em}%
}

% --- Selector de estilo -------------------------------------------------------
\newcommand{\estiloApuntes}{\let\problema\problemacaja\let\endproblema\endproblemacaja}
\newcommand{\estiloExamen}{\let\problema\ejercicio\let\endproblema\endejercicio}
\estiloApuntes  % por defecto, todo el proyecto original queda igual que antes

% ==============================================================================
% ENTORNOS DE EXÁMENES (del proyecto nuevo, sin cambios)
% ==============================================================================

% --- Entorno Resolución (separado por línea horizontal sencilla) --------------
\newenvironment{resolucion}{%
  \par\addvspace{0.4em}%
  \noindent\rule{\linewidth}{0.6pt}\par\nobreak
  \vspace{0.3em}%
  \noindent{\large\bfseries Resolución Paso a Paso:\par}\nobreak
  \vspace{0.5em}\noindent\ignorespaces
}{%
  \par\addvspace{0.4em}%
  \noindent\rule{\linewidth}{0.6pt}\par
  \addvspace{1.2em}%
}

% --- Entorno Mecanizado (estructura metodológica fija) ------------------------
\newenvironment{mecanizado}{%
  \par\addvspace{0.4em}%
  \small
  \noindent\textbf{Estructura de resolución}\par\nobreak
  \vspace{0.2em}%
  \itshape\noindent\ignorespaces
}{%
  \par\addvspace{0.6em}%
}

% --- Hueco de diagrama con descripción para IA --------------------------------
\newcommand{\huecodiagrama}[2]{%
  \par\addvspace{0.6em}%
  \begin{center}
    \framebox[\linewidth][c]{%
      \begin{minipage}{0.95\linewidth}
        \vspace{0.4em}
        \noindent\textbf{[DIAGRAMA REQUERIDO: #1]}\\[0.3em]
        \small\noindent\textbf{Descripción técnica para el generador gráfico:} #2\par
        \vspace{0.4em}
      \end{minipage}%
    }
  \end{center}
  \par\addvspace{0.6em}%
}